% !TeX root = ../bd-screen.tex

\chapter{引言}

模态逻辑是经典逻辑的扩展，通过在公式上附加算子 $\Box$（“方框”）和 $\Diamond$（“菱形”）而构成。
直观地说，$\Box$ 可以读作“必然”，$\Diamond$ 可以读作“可能”，因此 $\Box p$ 表示“$p$ 必然为真”，
而 $\Diamond p$ 表示“$p$ 可能为真”。由于必然性与可能性是根本的形而上学概念，模态逻辑显然具有重大的哲学意义。
它可以形式化诸如“$\Box p \lif p$”（若 $p$ 是必然的，则 $p$ 为真）或“$\Diamond p \lif \Box\Diamond p$”
（若 $p$ 是可能的，则它必然可能）这样的形而上学原则。

算子 $\Box$ 和 $\Diamond$ 是\emph{内涵的}。
这意味着 $\Box !A$ 或 $\Diamond !A$ 是否成立，并不只取决于 $!A$ 成立与否。
非内涵的算子是\emph{外延的}。否定是外延的：$\lnot !A$ 成立当且仅当 $!A$ 不成立；
所以 $\lnot !A$ 是否成立仅取决于 $!A$ 是否成立。
$\Box$ 和 $\Diamond$ 则不是这样：$\Box !A$ 或 $\Diamond !A$ 是否成立，还依赖于~$!A$ 的意义。
普通的真值函数语义足以处理外延算子，但像 $\Box$ 和 $\Diamond$ 这样的内涵算子则需要一种不同的语义。
本书中占据核心地位的一种这样的语义是关系语义（也称为可能世界语义或 Kripke（克里普克）语义）。

对于将 $\Box$ 解释为“必然”的逻辑而言，这种语义相对简单：解释~$\mModel{M}$ 不是给命题变元指派真值，
而是给它们指派一组“世界”——直观上，也就是 $p$ 在其中被解释为真的那些世界~$w$。
基于这样的解释，我们可以定义一个满足关系。
该满足关系的定义规定，$\Box !A$ 在世界~$w$ 处被满足，当且仅当 $!A$ 在\emph{所有}世界处都被满足：$\mSat{M}{\Box !A}[w]$当且仅当
$\mSat{M}{!A}[v]$ 对所有世界~$v$ 成立。这对应了 Leibniz（莱布尼茨）的观点：
必然为真的就是在所有可能世界中都为真的东西。

“必然”并非解释 $\Box$ 算子的唯一方式，但它是标准方式——“必然”和“可能”是所谓的\emph{真势}模态。
其他解释将 $\Box$ 读作“（某个人~$A$）知道”，或“某个人 $A$ 相信”，或“应当如此”，或“将来总是为真”。
这些分别是认知、信念、道义和时态模态。对 $\Box$ 的不同解释将使不同的公式成为逻辑真理，
并判定不同的推理为有效。例如，一切必然的与一切已知的都为真，因此 $\Box !A \lif !A$
在真势和认知解释下是逻辑真理。相反，并非一切被相信的或一切应当如此的事情实际上都为真，
所以 $\Box !A \lif !A$ 在信念或道义解释下不是逻辑真理。

为了处理模态算子的不同解释，语义通过在世界之间引入一个关系（即可达关系）而得到扩展。
于是 $\mSat{M}{\Box !A}[w]$当且仅当 $\mSat{M}{!A}[v]$对所有从~$w$ 可达的世界~$v$ 成立。
由此产生的语义非常灵活且强大，其基本思想可用于为基于其他内涵算子的逻辑提供语义解释。
其中之一是直觉主义逻辑，这是一种基于 L. E. J. Brouwer（布劳威尔）的构造主义数学分支的构造性逻辑。
直觉主义逻辑因此具有哲学上的趣味——它在数学的构造性论述中扮演重要角色——但二十世纪有影响力的
英国哲学家 Michael Dummett（达米特）也曾提出，它是一种优于经典逻辑的逻辑。
关系模型的另一个应用是作为虚拟条件句或反事实条件句的语义，这种方法由 Robert Stalnaker（斯塔纳克）和 David K. Lewis（刘易斯）开创。

本书旨在介绍内涵逻辑的语法、语义与证明理论。
目前它只处理命题逻辑，不过未来的版本也会涉及谓词逻辑。
全书内容分为三个部分：第一部分讨论正规模态逻辑，
即带有算子 $\Box$ 和~$\Diamond$ 的逻辑。
我们讨论它们的语法、关系模型以及基于这些模型的语义概念（如有效性与后承），
还有推导系统（包括公理系统和分析表）。
我们将建立关于这些逻辑的一些基本结果，
例如所考虑的推导系统的可靠性与完备性，
并讨论一些模型论构造，如过滤。
第二部分讨论直觉主义逻辑。
在此我们讨论自然演绎与公理推导、关系语义与拓扑语义，
以及推导系统的可靠性与完备性。
第三部分讨论反事实条件句的 Lewis-Stalnaker 语义。
附录则讨论了对于关系语义至关重要的集合论与关系论的一些概念和结果，
并回顾了经典命题逻辑的语法、语义和推导理论。